% Post-CP19 candidate; this extends only the theta jet geometry.
\subsection{Fixed-point jet finiteness without ordinarity}
\label{subsec:nonordinary-ppav-fixed-jets-candidate}
\status{Candidate; complete elementary proof written, independent audit pending.}

\begin{lemma}[The doubled image of a smooth principal theta curve]
\label{lem:all-theta-double-image-candidate}
Let $(A,\Theta)$ be a principally polarized abelian surface
over an algebraically closed field of characteristic two,
with smooth theta divisor. Then the finite map
\[
                         [2]|_\Theta:\Theta\longrightarrow[2]\Theta
\]
has inseparable degree two and separable degree one. In
particular its reduced image is numerically equivalent to
$2\Theta$.
\end{lemma}
\begin{proof}
The smooth theta curve has genus two and identifies $A$
with its Jacobian by Lemma~\ref{lem:ppav-surface-dichotomy}.
Write $[2]=V\circ F_{A/k}$, with
$V:A^{(2)}\to A$ the Verschiebung.
The differential $dV$ has the rank of the Cartier operator
on the two-dimensional space of canonical forms of $\Theta$.
Here is the relevant type identification: on the Jacobian,
$V$ is pullback of line bundles by the relative Frobenius
of the curve; its tangent map is the linearized Frobenius
on $H^1(\mathcal O_\Theta)$. Serre duality identifies its
dual with relative Cartier on $H^0(\omega_\Theta)$.
This is the same Picard/Frobenius and finite-duality
identification used in the preceding rank and covector
proofs.

This rank cannot be zero. Choose a separating rational
parameter $x$ on the smooth curve. In characteristic two
every rational differential killed by Cartier has the
form $b^2\,dx$: write its coefficient in the two-basis
$1,x$ over the field of squares and apply Cartier.
If both independent canonical forms were killed, their
ratio would be a square, say $z^2$. But their ratio
defines the canonical degree-two map to $\mathbf P^1$.
The equality
\[
 [k(\Theta):k(z^2)]
       =[k(\Theta):k(z)]\,[k(z):k(z^2)]=2
\]
would give $k(\Theta)=k(z)$, making the genus-two curve
rational. This is impossible. Thus $dV$ has rank one or
two.

If its rank is two, the restriction of $V$ to
$\Theta^{(2)}$ is separable. Suppose instead its rank is
one and that this curve restriction has zero differential.
The invariant line $\ker(dV)$ is then tangent to
$\Theta^{(2)}$ everywhere. Choose a nonzero invariant
vector field $\delta$ spanning it. Locally, if $f=0$
defines that smooth Cartier divisor, tangency says
$\delta(f)\in(f)$. First-order translation by
$\varepsilon\delta$, over $k[\varepsilon]/(\varepsilon^2)$,
therefore preserves the divisor: it sends its equation to
$f+\varepsilon\delta(f)=f(1+\varepsilon a)$.
It consequently preserves its line bundle. This places
the nonzero tangent vector $\delta$ in the scheme kernel
of the principal polarization
$\lambda_{\mathcal O(\Theta^{(2)})}$, which is an
isomorphism. The contradiction proves that
$V|_{\Theta^{(2)}}$ is separable also in rank one.
No integration of a vector field or reducedness assumption
on the kernel of $V$ is used here.

Its separable degree is one. Otherwise a generic fibre
would contain two distinct points differing by a nonzero
element of the finite set $\ker(V)(k)$. One such difference
would occur on a dense subset of the irreducible curve,
forcing a nonzero translation to preserve
$\Theta^{(2)}$. Its translation stabilizer is contained
scheme-theoretically in the kernel of its principal
polarization and is trivial. This is again impossible.

Relative Frobenius restricts on the smooth theta curve
with degree two. Thus $[2]|_\Theta$ has the asserted
degrees. By the projection formula,
$[2]_*\Theta\equiv4\Theta$ on the surface, since
$[2]^*$ multiplies divisor classes numerically by four.
Its generic curve degree is two, so the reduced image
satisfies $[2]\Theta\equiv2\Theta$.
\end{proof}

\begin{corollary}[The two-jet exceptional set has finite fixed fibres]
\label{cor:all-theta-fixed-point-jets-candidate}
In Theorem~\ref{thm:ppav-three-theta-jets}, the fixed-point
finiteness assertion holds for every smooth-theta abelian
surface in characteristic two, whether ordinary or not.
Explicitly, for each fixed $P\in A(k)$ the set
\[
 R_P=\{Q:2P+Q\in\Theta,\ P+2Q\in\Theta\}
\]
is finite. Consequently the first jets at $P,Q$ have full
rank six for all but finitely many $Q$ on every prescribed
integral ample curve.
\end{corollary}
\begin{proof}
If $R_P$ were infinite, its first condition would make it
dense in the irreducible curve $\Theta-2P$. Its second
condition would force
\[
                          [2]\Theta=\Theta+3P.
\]
The left-hand reduced curve has numerical class $2\Theta$
by the lemma, while the right-hand curve has class
$\Theta$, a contradiction. The rank-six assertion is
then exactly part two of the cited two-jet theorem.
\end{proof}

\begin{remark}[Scope]
\label{rem:all-theta-fixed-point-jets-scope-candidate}
This removes only the ordinarity restriction from the
fixed-point exceptional-jet argument. It asserts neither
relative Hasse-divisor irreducibility nor full monodromy
for a fixed nonproduct nonordinary surface. In particular
one cannot pad an odd sharp norm curve on a nonordinary
surface with another ample component and retain its sharp
rank: the positive component defects add on a nodal union.
The elementary first-order stabilizer proof above avoids
any unsupported extension of an ordinary Hecke or Hasse
argument.
\end{remark}
